\section{Calibration-dimension details}
\label{app:ccdim}

This appendix records the exact geometric objects used in
Theorem~\ref{thm:main}. Let $L\in\R_+^{n\times k}$ be a finite loss matrix,
with outcome distributions in $\Delta_n$. For report $t$, write
\[
 Q_t^L=\{q\in\Delta_n:q^\top L_{\cdot,t}
            \leq q^\top L_{\cdot,t'}\text{ for every }t'\in[k]\}.
 \tag{27}\label{eq:triggerapp}
\]
For a convex set $Q\subseteq\R^n$ and $p\in Q$, its cone of feasible
directions is
\[
 F_Q(p)=\{v:\text{there is }\epsilon_0>0\text{ such that }
 p+\epsilon v\in Q\text{ for every }\epsilon\in(0,\epsilon_0)\}.
\]
The feasible subspace is $F_Q(p)\cap(-F_Q(p))$, and its dimension is
$\mu_Q(p)$. Theorem 16 of \citet{ramaswamy2016ccdim} gives
\[
 \CCdim(L)\geq\|p\|_0-\mu_{Q_t^L}(p)-1
 \quad\text{for every }p\in Q_t^L.
 \tag{28}\label{eq:ramatheorem}
\]
It applies to the calibration definition used in their paper and does not
require the conditional surrogate risk to attain its infimum.

For completeness, their affine upper bound states
\[
 \CCdim(L)\leq\affdim(L),
 \tag{29}\label{eq:affupper}
\]
where $\affdim(L)$ is the affine dimension of the set of loss columns. One
way to view the construction is to choose affine coordinates for the loss
columns, estimate the corresponding conditional expectations with a convex
quadratic loss, and decode by minimizing the reconstructed target risk.

\section{Proof details for the rank transition}
\label{app:rank}

We prove the positive-$\lambda$ part of Theorem~\ref{thm:rank} with all
interchanges explicit. Let $\mathcal Y_+=2^{[s]}\setminus\{\varnothing\}$.
For $A,B\in\mathcal Y_+$, abbreviate
\[
 x=|A\cap B|,\qquad a=|A|,\qquad b=|B|.
\]
Because $0\leq\lambda<2$ and $x\leq\min(a,b)$, the denominator in
\eqref{eq:family} is positive. The elementary identity
\[
 \frac{x}{a+b-\lambda x}
 =\int_0^\infty xe^{-t(a+b-\lambda x)}dt
 \tag{30}\label{eq:laplaceapp}
\]
gives \eqref{eq:integral}.

For incidence vectors $\chi_A,\chi_B\in\{0,1\}^s$,
\[
 x^r=\langle\chi_A^{\otimes r},\chi_B^{\otimes r}\rangle.
 \tag{31}\label{eq:tensoridentity}
\]
Moreover,
\[
 xe^{\lambda tx}
 =\sum_{r=1}^\infty
 \frac{(\lambda t)^{r-1}}{(r-1)!}x^r.
 \tag{32}\label{eq:seriesapp}
\]
All coefficients in this series are nonnegative. Since the set family is
finite, Tonelli's theorem applies after separating the positive and negative
parts of the finite coefficient vector, or equivalently after first expanding
the finite quadratic form. Equations~\eqref{eq:laplaceapp}--\eqref{eq:seriesapp}
yield
\begin{align}
 \sum_{A,B\in\mathcal Y_+}c_Ac_BS_\lambda(A,B)
 &=(2-\lambda)\int_0^\infty
 \sum_{r=1}^\infty\frac{(\lambda t)^{r-1}}{(r-1)!}
 \left\langle
 \sum_Ac_Ae^{-t|A|}\chi_A^{\otimes r},
 \sum_Bc_Be^{-t|B|}\chi_B^{\otimes r}
 \right\rangle dt\\
 &=(2-\lambda)\int_0^\infty
 \sum_{r=1}^\infty\frac{(\lambda t)^{r-1}}{(r-1)!}
 \left\|\sum_Ac_Ae^{-t|A|}\chi_A^{\otimes r}\right\|_2^2dt.
 \tag{33}\label{eq:quadraticapp}
\end{align}
The integrand is continuous and nonnegative. If the integral is zero, the
integrand is zero for every $t>0$. Since $\lambda>0$, every tensor norm in
the sum is zero.

Fix a nonempty $T=\{i_1,\ldots,i_r\}\subseteq[s]$. The
$(i_1,\ldots,i_r)$ coordinate of the $r$th tensor sum is
\[
 \sum_{A\in\mathcal Y_+}c_Ae^{-t|A|}
 \prod_{j=1}^r\chi_A(i_j)
 =\sum_{A\supseteq T}c_Ae^{-t|A|}.
 \tag{34}\label{eq:zetacoord}
\]
Fix any $t>0$ and put $d_A=c_Ae^{-t|A|}$. Equation~\eqref{eq:zetacoord}
says that the upper-zeta transform $T\mapsto\sum_{A\supseteq T}d_A$ is zero
on the poset of nonempty subsets. Ordered by decreasing cardinality, its
matrix is triangular with ones on the diagonal. Hence $d_A=0$ for every
$A$, and therefore $c_A=0$. The nonempty block is strictly positive
definite. The empty set has score zero against every nonempty set and score
one against itself, so it contributes a separate positive $1\times1$ block.

At $\lambda=0$, only the $r=1$ term remains in
\eqref{eq:seriesapp}. The variation of $e^{-t|A|}$ with cardinality allows
more than $s$ rank, but not full rank. \citet{zhang2026exact} prove the exact
value $s^2-s+2$, including the empty-set block.

Finally, if a square matrix $K$ of size $N$ is nonsingular, its $N$ columns
are linearly independent and thus affinely independent. The affine map
$k_B\mapsto\1-k_B$ preserves affine independence, so the loss columns of
$\1\1^\top-K$ have affine dimension $N-1$. This proves
\eqref{eq:affdim}.

\section{The factorial tie identity}
\label{app:tie}

We give a full proof of Lemma~\ref{lem:tie}. Recall that $O$ contains the
$n=s-1$ optional labels. For $C\subseteq O$, define $W(C)$ by
\eqref{eq:W}. We first show that $W(C)$ is invariant under adding an element.

Fix $j\in O\setminus C$. Pair the outcome index $D=E$ with
$D=E\cup\{j\}$ for each $E\subseteq O\setminus\{j\}$. Put
\[
 a=|E\cap C|,\qquad b=|E\setminus C|,\qquad c=|C|.
\]
For $D=E$, adding $j$ to the report leaves the intersection unchanged and
increases the union size by one. Its contribution to
$W(C\cup\{j\})-W(C)$ is
\[
 -\frac{a+1}{(a+b)!(c+b+1)(c+b+2)}.
 \tag{35}\label{eq:nojinD}
\]
For $D=E\cup\{j\}$, adding $j$ to the report leaves the union unchanged and
increases the intersection by one. Its contribution is
\[
 \frac1{(a+b+1)!(c+b+2)}.
 \tag{36}\label{eq:jinD}
\]
Their sum is \eqref{eq:increment}.

For a fixed $b$, the number of choices of $E$ having a given $a$ is
$\binom ca\binom{n-c-1}{b}$. The second binomial coefficient and the factor
$1/(c+b+2)$ do not depend on $a$. It remains to prove
\eqref{eq:identity}. Multiply its left-hand side by $c+b+1$ and use
$c+b+1=(c-a)+(a+b+1)$:
\begin{align}
 (c+b+1)\sum_{a=0}^c\binom ca\frac1{(a+b+1)!}
 &=\sum_{a=0}^c\binom ca\frac1{(a+b)!}
   +\sum_{a=0}^c(c-a)\binom ca\frac1{(a+b+1)!}\\
 &=\sum_{a=0}^c\binom ca\frac1{(a+b)!}
   +c\sum_{a=0}^{c-1}\binom{c-1}{a}\frac1{(a+b+1)!}\\
 &=\sum_{a=0}^c\binom ca\frac1{(a+b)!}
   +\sum_{a=0}^ca\binom ca\frac1{(a+b)!}\\
 &=\sum_{a=0}^c(a+1)\binom ca\frac1{(a+b)!}.
 \tag{37}\label{eq:identityproof}
\end{align}
In the third line we shifted the index and used
$a\binom ca=c\binom{c-1}{a-1}$. Hence every fixed-$b$ sum of paired
increments is zero, and $W(C\cup\{j\})=W(C)$.

Repeatedly removing elements from $C$ gives $W(C)=W(\varnothing)$. Directly,
\[
 W(\varnothing)
 =\sum_{D\subseteq O}\frac1{|D|!(1+|D|)}
 =\sum_{d=0}^n\binom nd\frac1{(d+1)!},
 \tag{38}\label{eq:wempty}
\]
which proves \eqref{eq:constant}. After division by $Z_n$, all reports in
$\mathcal U$ therefore tie.

Now take any report $B$ with $\star\notin B$ and put
$B^+=B\cup\{\star\}$. Every outcome $A$ in the support of $p$ contains
$\star$. Therefore
\[
 |B^+\cup A|=|B\cup A|,
 \qquad
 |B^+\cap A|=|B\cap A|+1.
\]
The Jaccard score of $B^+$ is strictly larger for every such $A$. Since all
support probabilities are positive, $B$ is strictly suboptimal. Thus exactly
$\mathcal U$ is Bayes optimal.

\section{Feasible-subspace calculation}
\label{app:feasible}

We expand the geometric step in Theorem~\ref{thm:main}. Let
$K_{\mathcal U,\mathcal U}$ be the $r\times r$ principal Jaccard block, where
$r=2^{s-1}$. It is positive definite by Theorem~\ref{thm:rank}. Fix
$B_0\in\mathcal U$ and form the matrix whose $r-1$ columns are
\[
 d_B=K_{\mathcal U,B}-K_{\mathcal U,B_0},
 \qquad B\in\mathcal U\setminus\{B_0\}.
 \tag{39}\label{eq:dBapp}
\]
These columns are independent. Indeed, if
$\sum_{B\ne B_0}\alpha_Bd_B=0$, define coefficients $\beta_B=\alpha_B$ for
$B\ne B_0$ and $\beta_{B_0}=-\sum_{B\ne B_0}\alpha_B$. Then
$K_{\mathcal U,\mathcal U}\beta=0$. Nonsingularity gives $\beta=0$, hence
every $\alpha_B=0$. Thus the matrix $D$ of differences has rank $r-1$.

The tie identity says $D^\top p_{\mathcal U}=0$. Since
$\dim\ker(D^\top)=1$ and $p_{\mathcal U}\ne0$,
\[
 \ker(D^\top)=\operatorname{span}(p_{\mathcal U}).
 \tag{40}\label{eq:kernelapp}
\]

Let $v\in F_{Q_{B_0}}(p)\cap(-F_{Q_{B_0}}(p))$. If $A\notin\mathcal U$,
then $p_A=0$. Feasibility of both $v$ and $-v$ in the nonnegative orthant
forces $v_A=0$. For every $B\in\mathcal U$, the target risks of $B$ and
$B_0$ tie at $p$. The corresponding inequality in $Q_{B_0}$ is active.
Feasibility in both directions forces its directional derivative to be zero:
\[
 v_{\mathcal U}^\top
 (L^{\rm Jac}_{\mathcal U,B}-L^{\rm Jac}_{\mathcal U,B_0})=0.
 \tag{41}\label{eq:activeapp}
\]
Since loss differences are negative score differences,
\eqref{eq:activeapp} is $D^\top v_{\mathcal U}=0$. By
\eqref{eq:kernelapp}, $v_{\mathcal U}=\alpha p_{\mathcal U}$. Finally,
two-sided feasibility in the affine hull of the simplex requires
$\1^\top v=0$. Because $\1^\top p=1$, this implies $\alpha=0$. The feasible
subspace is $\{0\}$ and has dimension zero.

Notice that constraints against reports outside $\mathcal U$ were not needed.
They are strict at $p$ by Lemma~\ref{lem:tie} and therefore impose no
additional equality on a two-sided feasible direction.

\section{Exact analysis of the logistic counterexample}
\label{app:counterexample}

For the set order
$\varnothing,\{1\},\{2\},\{1,2\}$, the Jaccard score matrix is
\[
 K=
 \begin{pmatrix}
 1&0&0&0\\
 0&1&0&1/2\\
 0&0&1&1/2\\
 0&1/2&1/2&1
 \end{pmatrix}.
 \tag{42}\label{eq:Ktwo}
\]
Multiplying \eqref{eq:counterp} by this matrix gives
\eqref{eq:countera}. The unique maximum $3/4$ occurs at report $\{1\}$.

We next derive the conditional surrogate risk. Let
$r_B=\operatorname{score}_h(B)$. Each term in \eqref{eq:maoloss} satisfies
\[
 \log\sum_Ce^{r_C-r_B}
 =\log\sum_Ce^{r_C}-r_B=-\log s_B(h).
 \tag{43}\label{eq:logsoftmaxapp}
\]
Taking expectation over $Y$ and interchanging two finite sums gives
\[
 \E_p\widetilde L(h,Y)
 =-\sum_B\E_p[\Jac(B,Y)]\log s_B(h)
 =-\sum_Ba_B\log s_B(h).
 \tag{44}\label{eq:riskapp}
\]

Because $r_B=\sum_i y_i^Bh_i$, $s_B(h)$ is the product of two Bernoulli
probabilities. If $q_B=a_B/A$ with $A=\sum_Ba_B=33/20$, then
\eqref{eq:riskapp} is $A$ times the cross entropy from $q$ to this product
family. Writing $\pi_i=\Pr_{s(h)}(i\in B)$, the part depending on $h$ is
\[
 A\sum_{i=1}^2[-m_i\log\pi_i-(1-m_i)\log(1-\pi_i)],
 \tag{45}\label{eq:binaryce}
\]
up to a constant independent of $h$. Each binary cross entropy is strictly
convex in its logit and uniquely minimized at $\pi_i=m_i$. Direct arithmetic
gives
\[
 A=\frac{33}{20},\qquad
 m_1=\frac{3/4+13/20}{33/20}=\frac{28}{33},\qquad
 m_2=\frac{1/4+13/20}{33/20}=\frac6{11}.
 \tag{46}\label{eq:marginalarithmetic}
\]
Under the $\{-1,1\}$ score encoding,
$\pi_i=e^{h_i}/(e^{h_i}+e^{-h_i})$. Solving for $h_i$ yields
\eqref{eq:hstar}. Both coordinates are positive, so the sign link returns
$\{1,2\}$. Its target loss exceeds the Bayes loss by
\[
 \left(1-\frac{13}{20}\right)-\left(1-\frac34\right)=\frac1{10}.
 \tag{47}\label{eq:regretapp}
\]
The conditional surrogate risk is minimized exactly, so its regret is zero.

For reference, the product constraint \eqref{eq:productconstraint} follows
immediately from the four unnormalized weights
\[
 e^{-h_1-h_2},\quad e^{h_1-h_2},\quad
 e^{-h_1+h_2},\quad e^{h_1+h_2}.
\]
Both cross-products equal one before applying the common softmax normalizer.

\section{Finite exact checks}
\label{app:checks}

The proof is symbolic, but small cases provide useful falsification checks.
For $s=2,3,4,5$, exact rational calculations for the witness
\eqref{eq:witness} give
\[
\begin{array}{c|cccc}
s & |\mathcal U| & \text{tied score} & \text{best score outside }\mathcal U
  & \operatorname{rank}(K_{\mathcal U,\mathcal U})\\ \hline
2&2&3/4&1/4&2\\
3&4&13/21&2/7&4\\
4&8&73/136&39/136&8\\
5&16&501/1045&292/1045&16
\end{array}
\tag{48}\label{eq:smallchecks}
\]

The calculation enumerates all $2^s$ outcomes and reports using rational
arithmetic, checks equality of the scores over $\mathcal U$, strict
separation outside $\mathcal U$, and exact rank of the rational principal
matrix. The two-label certificate separately checks
\[
 a=(0,3/4,1/4,13/20),\quad
 (m_1,m_2)=(28/33,6/11),\quad
 \text{regret}=1/10.
\]
No numerical approximation is used in these checks.
